\documentclass[a4paper,12pt,twoside]{memoir}

% Castellano
\usepackage[spanish,es-tabla]{babel}
\selectlanguage{spanish}
\usepackage[utf8]{inputenc}
\usepackage[T1]{fontenc}
\usepackage{lmodern} % Scalable font
\usepackage{microtype}
\usepackage{placeins}

\RequirePackage{booktabs}
\RequirePackage[table]{xcolor}
\RequirePackage{xtab}
\RequirePackage{multirow}

% Links
\PassOptionsToPackage{hyphens}{url}\usepackage[colorlinks]{hyperref}
\hypersetup{
	allcolors = {red}
}

% Ecuaciones
\usepackage{amsmath}

% Rutas de fichero / paquete
\newcommand{\ruta}[1]{{\sffamily #1}}

% Párrafos
\nonzeroparskip

% Huérfanas y viudas
\widowpenalty100000
\clubpenalty100000

% Imágenes

% Comando para insertar una imagen en un lugar concreto.
% Los parámetros son:
% 1 --> Ruta absoluta/relativa de la figura
% 2 --> Texto a pie de figura
% 3 --> Tamaño en tanto por uno relativo al ancho de página
\usepackage{graphicx}
\newcommand{\imagen}[3]{
	\begin{figure}[!h]
		\centering
		\includegraphics[width=#3\textwidth]{#1}
		\caption{#2}\label{fig:#1}
	\end{figure}
	\FloatBarrier
}

% Comando para insertar una imagen sin posición.
% Los parámetros son:
% 1 --> Ruta absoluta/relativa de la figura
% 2 --> Texto a pie de figura
% 3 --> Tamaño en tanto por uno relativo al ancho de página
\newcommand{\imagenflotante}[3]{
	\begin{figure}
		\centering
		\includegraphics[width=#3\textwidth]{#1}
		\caption{#2}\label{fig:#1}
	\end{figure}
}

% El comando \figura nos permite insertar figuras comodamente, y utilizando
% siempre el mismo formato. Los parametros son:
% 1 --> Porcentaje del ancho de página que ocupará la figura (de 0 a 1)
% 2 --> Fichero de la imagen
% 3 --> Texto a pie de imagen
% 4 --> Etiqueta (label) para referencias
% 5 --> Opciones que queramos pasarle al \includegraphics
% 6 --> Opciones de posicionamiento a pasarle a \begin{figure}
\newcommand{\figuraConPosicion}[6]{%
  \setlength{\anchoFloat}{#1\textwidth}%
  \addtolength{\anchoFloat}{-4\fboxsep}%
  \setlength{\anchoFigura}{\anchoFloat}%
  \begin{figure}[#6]
    \begin{center}%
      \Ovalbox{%
        \begin{minipage}{\anchoFloat}%
          \begin{center}%
            \includegraphics[width=\anchoFigura,#5]{#2}%
            \caption{#3}%
            \label{#4}%
          \end{center}%
        \end{minipage}
      }%
    \end{center}%
  \end{figure}%
}

%
% Comando para incluir imágenes en formato apaisado (sin marco).
\newcommand{\figuraApaisadaSinMarco}[5]{%
  \begin{figure}%
    \begin{center}%
    \includegraphics[angle=90,height=#1\textheight,#5]{#2}%
    \caption{#3}%
    \label{#4}%
    \end{center}%
  \end{figure}%
}
% Para las tablas
\newcommand{\otoprule}{\midrule [\heavyrulewidth]}
%
% Nuevo comando para tablas pequeñas (menos de una página).
\newcommand{\tablaSmall}[5]{%
 \begin{table}
  \begin{center}
   \rowcolors {2}{gray!35}{}
   \begin{tabular}{#2}
    \toprule
    #4
    \otoprule
    #5
    \bottomrule
   \end{tabular}
   \caption{#1}
   \label{tabla:#3}
  \end{center}
 \end{table}
}

%
% Nuevo comando para tablas pequeñas (menos de una página).
\newcommand{\tablaSmallSinColores}[5]{%
 \begin{table}[H]
  \begin{center}
   \begin{tabular}{#2}
    \toprule
    #4
    \otoprule
    #5
    \bottomrule
   \end{tabular}
   \caption{#1}
   \label{tabla:#3}
  \end{center}
 \end{table}
}

\newcommand{\tablaApaisadaSmall}[5]{%
\begin{landscape}
  \begin{table}
   \begin{center}
    \rowcolors {2}{gray!35}{}
    \begin{tabular}{#2}
     \toprule
     #4
     \otoprule
     #5
     \bottomrule
    \end{tabular}
    \caption{#1}
    \label{tabla:#3}
   \end{center}
  \end{table}
\end{landscape}
}

%
% Nuevo comando para tablas grandes con cabecera y filas alternas coloreadas en gris.
\newcommand{\tabla}[6]{%
  \begin{center}
    \tablefirsthead{
      \toprule
      #5
      \otoprule
    }
    \tablehead{
      \multicolumn{#3}{l}{\small\sl continúa desde la página anterior}\\
      \toprule
      #5
      \otoprule
    }
    \tabletail{
      \hline
      \multicolumn{#3}{r}{\small\sl continúa en la página siguiente}\\
    }
    \tablelasttail{
      \hline
    }
    \bottomcaption{#1}
    \rowcolors {2}{gray!35}{}
    \begin{xtabular}{#2}
      #6
      \bottomrule
    \end{xtabular}
    \label{tabla:#4}
  \end{center}
}

%
% Nuevo comando para tablas grandes con cabecera.
\newcommand{\tablaSinColores}[6]{%
  \begin{center}
    \tablefirsthead{
      \toprule
      #5
      \otoprule
    }
    \tablehead{
      \multicolumn{#3}{l}{\small\sl continúa desde la página anterior}\\
      \toprule
      #5
      \otoprule
    }
    \tabletail{
      \hline
      \multicolumn{#3}{r}{\small\sl continúa en la página siguiente}\\
    }
    \tablelasttail{
      \hline
    }
    \bottomcaption{#1}
    \begin{xtabular}{#2}
      #6
      \bottomrule
    \end{xtabular}
    \label{tabla:#4}
  \end{center}
}

%
% Nuevo comando para tablas grandes sin cabecera.
\newcommand{\tablaSinCabecera}[5]{%
  \begin{center}
    \tablefirsthead{
      \toprule
    }
    \tablehead{
      \multicolumn{#3}{l}{\small\sl continúa desde la página anterior}\\
      \hline
    }
    \tabletail{
      \hline
      \multicolumn{#3}{r}{\small\sl continúa en la página siguiente}\\
    }
    \tablelasttail{
      \hline
    }
    \bottomcaption{#1}
  \begin{xtabular}{#2}
    #5
   \bottomrule
  \end{xtabular}
  \label{tabla:#4}
  \end{center}
}



\definecolor{cgoLight}{HTML}{EEEEEE}
\definecolor{cgoExtralight}{HTML}{FFFFFF}

%
% Nuevo comando para tablas grandes sin cabecera.
\newcommand{\tablaSinCabeceraConBandas}[5]{%
  \begin{center}
    \tablefirsthead{
      \toprule
    }
    \tablehead{
      \multicolumn{#3}{l}{\small\sl continúa desde la página anterior}\\
      \hline
    }
    \tabletail{
      \hline
      \multicolumn{#3}{r}{\small\sl continúa en la página siguiente}\\
    }
    \tablelasttail{
      \hline
    }
    \bottomcaption{#1}
    \rowcolors[]{1}{cgoExtralight}{cgoLight}

  \begin{xtabular}{#2}
    #5
   \bottomrule
  \end{xtabular}
  \label{tabla:#4}
  \end{center}
}



\graphicspath{ {./img/} }

% Capítulos
\chapterstyle{bianchi}
\newcommand{\capitulo}[2]{
	\setcounter{chapter}{#1}
	\setcounter{section}{0}
	\setcounter{figure}{0}
	\setcounter{table}{0}
	\chapter*{#2}
	\addcontentsline{toc}{chapter}{#2}
	\markboth{#2}{#2}
}

% Apéndices
\renewcommand{\appendixname}{Apéndice}
\renewcommand*\cftappendixname{\appendixname}

\newcommand{\apendice}[1]{
	%\renewcommand{\thechapter}{A}
	\chapter{#1}
}

\renewcommand*\cftappendixname{\appendixname\ }

% Formato de portada
\makeatletter
\usepackage{xcolor}
\newcommand{\tutor}[1]{\def\@tutor{#1}}
\newcommand{\course}[1]{\def\@course{#1}}
\definecolor{cpardoBox}{HTML}{E6E6FF}
\def\maketitle{
  \null
  \thispagestyle{empty}
  % Cabecera ----------------
\noindent\includegraphics[width=\textwidth]{cabecera}\vspace{1cm}%
  \vfill
  % Título proyecto y escudo informática ----------------
  \colorbox{cpardoBox}{%
    \begin{minipage}{.8\textwidth}
      \vspace{.5cm}\Large
      \begin{center}
      \textbf{TFG del Grado en Ingeniería Informática}\vspace{.6cm}\\
      \textbf{\LARGE\@title{}}
      \end{center}
      \vspace{.2cm}
    \end{minipage}

  }%
  \hfill\begin{minipage}{.20\textwidth}
    \includegraphics[width=\textwidth]{escudoInfor}
  \end{minipage}
  \vfill
  % Datos de alumno, curso y tutores ------------------
  \begin{center}%
  {%
    \noindent\LARGE
    Presentado por \@author{}\\ 
    en Universidad de Burgos --- \@date{}\\
    Tutor: \@tutor{}\\
  }%
  \end{center}%
  \null
  \cleardoublepage
  }
\makeatother

\newcommand{\nombre}{Nombre del alumno} %%% cambio de comando

% Datos de portada
\title{título del TFG}
\author{\nombre}
\tutor{nombre tutor}
\date{\today}

\begin{document}

\maketitle


\newpage\null\thispagestyle{empty}\newpage


%%%%%%%%%%%%%%%%%%%%%%%%%%%%%%%%%%%%%%%%%%%%%%%%%%%%%%%%%%%%%%%%%%%%%%%%%%%%%%%%%%%%%%%%
\thispagestyle{empty}


\noindent\includegraphics[width=\textwidth]{cabecera}\vspace{1cm}

\noindent D. nombre tutor, profesor del departamento de nombre departamento, área de nombre área.

\noindent Expone:

\noindent Que el alumno D. \nombre, con DNI dni, ha realizado el Trabajo final de Grado en Ingeniería Informática titulado título de TFG. 

\noindent Y que dicho trabajo ha sido realizado por el alumno bajo la dirección del que suscribe, en virtud de lo cual se autoriza su presentación y defensa.

\begin{center} %\large
En Burgos, {\large \today}
\end{center}

\vfill\vfill\vfill

% Author and supervisor
\begin{minipage}{0.45\textwidth}
\begin{flushleft} %\large
Vº. Bº. del Tutor:\\[2cm]
D. nombre tutor
\end{flushleft}
\end{minipage}
\hfill
\begin{minipage}{0.45\textwidth}
\begin{flushleft} %\large
Vº. Bº. del co-tutor:\\[2cm]
D. nombre co-tutor
\end{flushleft}
\end{minipage}
\hfill

\vfill

% para casos con solo un tutor comentar lo anterior
% y descomentar lo siguiente
%Vº. Bº. del Tutor:\\[2cm]
%D. nombre tutor


\newpage\null\thispagestyle{empty}\newpage




\frontmatter

% Abstract en castellano
\renewcommand*\abstractname{Resumen}
\begin{abstract}
En este primer apartado se hace una \textbf{breve} presentación del tema que se aborda en el proyecto.
\end{abstract}

\renewcommand*\abstractname{Descriptores}
\begin{abstract}
Palabras separadas por comas que identifiquen el contenido del proyecto Ej: servidor web, buscador de vuelos, android \ldots
\end{abstract}

\clearpage

% Abstract en inglés
\renewcommand*\abstractname{Abstract}
\begin{abstract}
A \textbf{brief} presentation of the topic addressed in the project.
\end{abstract}

\renewcommand*\abstractname{Keywords}
\begin{abstract}
keywords separated by commas.
\end{abstract}

\clearpage

% Indices
\tableofcontents

\clearpage

\listoffigures

\clearpage

\listoftables
\clearpage

\mainmatter
\capitulo{1}{Introducción}

Green In House es una maceta educativa móvil diseñada para fomentar el aprendizaje de los niños y promover su conexión con la naturaleza. Es un producto que combina la tecnología y la educación ambiental de una manera interactiva y práctica.

Se trata de una maceta fabricada con materiales ecológicos, como el PETG reciclado, que se imprime utilizando tecnología de impresión 3D. Green In House no solo sirve como un macetero para plantas, sino que incorpora sensores y una aplicación móvil interactiva para brindar una experiencia de aprendizaje completa.

El objetivo de Green In House es mantener un registro temporal de los factores clave que influyen en el crecimiento de las plantas y enseñar a los niños sobre el cuidado de las mismas, concienciandolos sobre la importancia de la naturaleza. A través de la maceta educativa, los niños aprenden a asumir responsabilidades, como regar las plantas y proporcionarles un entorno adecuado.

La aplicación móvil complementaria permite a los niños acceder a información detallada sobre las plantas, recibir y generar consejos de cuidado y realizar un seguimiento del crecimiento de sus plantas. Además, la aplicación fomenta la interacción social al permitir a los niños compartir sus experiencias y consejos con otros usuarios.

Green In House también destaca por su enfoque en aspectos clave para el crecimiento de las plantas, como la temperatura, la humedad y la cantidad de luz. Los sensores incorporados en la maceta educativa permiten a los niños monitorizar y ajustar estos factores para garantizar un entorno óptimo para el desarrollo de las plantas.

En resumen, Green In House es una maceta educativa móvil que combina tecnología, educación ambiental y aprendizaje práctico. Proporciona a los niños la oportunidad de cultivar sus propias plantas, aprender sobre el cuidado de la naturaleza y adquirir habilidades valiosas para el cuidado del medio ambiente.





#TODO

Estructura de la memoria....


De que va el proyecto, que resolver (sin entrar en mucho detalle)

Descripción del contenido del trabajo y del estrucutra de la memoria y del resto de materiales entregados.

\capitulo{2}{Objetivos del proyecto}

Green In House es un proyecto que integra tecnología, ecología y educación de una manera práctica e interactiva, brindando a los niños la oportunidad de aprender sobre la importancia de cuidar el medio ambiente y asumir responsabilidades. Por lo tanto, Green In House no es solo una maceta, sino que es un instrumento de aprendizaje que promueve la comprensión de conceptos científicos como la fotosíntesis, la biología de las plantas y la interacción con el medio ambiente. Los sensores incorporados en la maceta ayudan a los niños a entender cómo diferentes factores como la luz, la humedad y la temperatura afectan el crecimiento de las plantas.

Además, la integración con una aplicación móvil proporciona una experiencia de aprendizaje más rica. La aplicación ofrece a los niños acceso a información detallada sobre las plantas, consejos de cuidado y la capacidad de registrar observaciones y hacer un seguimiento del crecimiento de sus plantas. Debido a esto es una herramienta educativa muy valiosa, ya que puede ayudar al aprendizaje infantil de manera significativa, vivencial y con aporte al aprendizaje académico de varias formas, a la vez que desarrolla nuevas habilidades trasversales gracias a la motivación:

\textbf{Aprendizaje significativo:} Green In House proporciona a los niños una experiencia práctica y tangible de cuidado de plantas, lo que les permite establecer conexiones significativas entre los conceptos teóricos y la realidad. Al participar activamente en el proceso de crecimiento de las plantas, los niños pueden comprender mejor los conceptos relacionados con la biología, la ecología, la fotosíntesis y otros temas relacionados.

\textbf{Aprendizaje vivencial:} Green In House ofrece una experiencia de aprendizaje práctica y tangible, donde los niños pueden interactuar directamente con las plantas y experimentar las consecuencias de sus acciones. Aprenden observando el crecimiento de las plantas, experimentando con diferentes condiciones ambientales y tomando decisiones sobre el cuidado de las mismas. Esta experiencia vivencial refuerza el aprendizaje al involucrar múltiples sentidos y brindar una experiencia multisensorial.

\textbf{Aporte al aprendizaje académico:} Green In House se puede integrar en el currículo académico, complementando y enriqueciendo las lecciones tradicionales. Las actividades relacionadas con las plantas pueden ser utilizadas para abordar diferentes áreas del conocimiento, como ciencias naturales, matemáticas (al medir el crecimiento de las plantas), lenguaje (a través de la comunicación y expresión de ideas sobre las plantas), arte (mediante la creación de proyectos artísticos inspirados en la naturaleza) y más. Esto amplía el alcance del aprendizaje académico y lo hace más contextualizado y relevante para los niños.

\textbf{Desarrollo de habilidades transversales:} Green In House promueve el desarrollo de habilidades transversales importantes, como la responsabilidad, el pensamiento crítico, la resolución de problemas, la toma de decisiones y la autonomía. Los niños aprenden a cuidar de las plantas, a tomar decisiones basadas en la observación y a solucionar problemas relacionados con el crecimiento de las mismas. Estas habilidades son transferibles a otras áreas de su vida y contribuyen a su desarrollo integral.

\textbf{Motivación y engagement:} El enfoque lúdico y práctico de Green In House captura el interés de los niños y los motiva a participar activamente en el aprendizaje. Al estar involucrados en el cuidado de las plantas, los niños se sienten más comprometidos y conectados con el proceso de aprendizaje, lo que favorece un mayor nivel de retención y comprensión de los contenidos.

Green In House, a pesar de estar basado en la tecnología e incorporar el uso de dispositivos móviles, está diseñado para promover la interacción con el mundo real y mejorar el desarrollo cognitivo de los niños. Algunas razones por las cuales este proyecto podría tener un impacto positivo en lugar de fomentar una dependencia negativa de la tecnología:

\textbf{Interacción con el mundo real}: Aunque Green In House usa una aplicación móvil, la mayoría de las acciones necesarias para el cuidado de la planta se realizan en el mundo físico. Los niños tienen que regar la planta, colocarla en un lugar donde reciba suficiente luz y monitorear los parámetros del ambiente físico. Esto significa que pasan un tiempo significativo interactuando con el mundo real, no solo con la tecnología.

\textbf{Consecuencias tangibles}: A diferencia de muchos juegos digitales, donde las acciones pueden no tener consecuencias reales, en Green In House las decisiones y acciones de los niños tienen resultados tangibles. Si olvidan regar su planta, verán cómo se marchita. Esto puede ayudarles a comprender que sus acciones tienen consecuencias en el mundo real.

\textbf{Fomenta la responsabilidad}: El cuidado regular de una planta puede enseñar a los niños sobre la responsabilidad y la gestión del tiempo. Tienen que recordar regar la planta y comprobar sus condiciones regularmente. Este es un hábito que puede transferirse a otras áreas de sus vidas.

\textbf{Uso de la tecnología con propósito}: Green In House no promueve el uso de la tecnología por el simple hecho de usarla. En cambio, la tecnología se utiliza como una herramienta para aprender y cuidar de algo vivo. Esto puede ayudar a los niños a entender que la tecnología tiene muchos usos útiles y beneficiosos, más allá del entretenimiento.

Por lo tanto, aunque Green In House integra la tecnología, lo hace de una manera que promueve la interacción positiva y la comprensión del mundo real. Al proporcionar una experiencia de aprendizaje práctica y consecuencias reales para las acciones, puede ayudar a los niños a desarrollar habilidades valiosas y a entender cómo sus acciones afectan al mundo que les rodea.

Green In House, a pesar de estar basado en la tecnología e incorporar el uso de dispositivos móviles, está diseñado para promover la interacción con el mundo real y mejorar el desarrollo cognitivo de los niños. Aquí hay algunas razones por las cuales este proyecto podría tener un impacto positivo en lugar de fomentar una dependencia negativa de la tecnología:

Interacción con el mundo real: Aunque Green In House usa una aplicación móvil, la mayoría de las acciones necesarias para el cuidado de la planta se realizan en el mundo físico. Los niños tienen que regar la planta, colocarla en un lugar donde reciba suficiente luz y monitorear los parámetros del ambiente físico. Esto significa que pasan un tiempo significativo interactuando con el mundo real, no solo con la tecnología.

Consecuencias tangibles: A diferencia de muchos juegos digitales, donde las acciones pueden no tener consecuencias reales, en Green In House las decisiones y acciones de los niños tienen resultados tangibles. Si olvidan regar su planta, verán cómo se marchita. Esto puede ayudarles a comprender que sus acciones tienen consecuencias en el mundo real.

Fomenta la responsabilidad: El cuidado regular de una planta puede enseñar a los niños sobre la responsabilidad y la gestión del tiempo. Tienen que recordar regar la planta y comprobar sus condiciones regularmente. Este es un hábito que puede transferirse a otras áreas de sus vidas.

Uso de la tecnología con propósito: Green In House no promueve el uso de la tecnología por el simple hecho de usarla. En cambio, la tecnología se utiliza como una herramienta para aprender y cuidar de algo vivo. Esto puede ayudar a los niños a entender que la tecnología tiene muchos usos útiles y beneficiosos, más allá del entretenimiento.

Por lo tanto, aunque Green In House integra la tecnología, lo hace de una manera que promueve la interacción positiva y la comprensión del mundo real. Al proporcionar una experiencia de aprendizaje práctica y consecuencias reales para las acciones, puede ayudar a los niños a desarrollar habilidades valiosas y a entender cómo sus acciones afectan al mundo que les rodea.


En resumen, Green In House contribuye al aprendizaje infantil significativo, vivencial y con aporte al aprendizaje académico al brindar experiencias prácticas y reales de cuidado de plantas, integrar conceptos teóricos en un contexto tangible, desarrollar habilidades transversales y promover la motivación y el engagement en el proceso de aprendizaje. Aunque pricipalmente está diseñado y pensado para niños, también puede ser útil para los adultos y hay varias segmentaciones de clientes que podrían encontrar este producto interesante:

Jardineros aficionados: Los adultos que disfrutan del jardín y quieren aprender más sobre el cuidado de las plantas podrían beneficiarse de las características interactivas y educativas de Green In House.

Amantes de la tecnología: Aquellos que están interesados en la tecnología, especialmente la tecnología que se integra con la vida cotidiana de formas nuevas e interesantes, pueden ver el valor en un producto como Green In House.

Profesores y educadores: Los profesores podrían utilizar Green In House como una herramienta educativa en las aulas para enseñar sobre ciencias naturales, ecología y biología, o para fomentar la responsabilidad y el cuidado del medio ambiente.

Padres: Los padres que quieren proporcionar experiencias educativas prácticas para sus hijos en casa podrían estar interesados en Green In House.

Personas que viven en espacios reducidos: Para aquellos que viven en apartamentos o casas sin jardín, Green In House ofrece una forma de tener y cuidar plantas en espacios interiores.

Empresas de bienestar en el trabajo: Las empresas que buscan mejorar el bienestar de sus empleados pueden estar interesadas en utilizar Green In House para animar a los empleados a tomar descansos activos, cuidar una planta y desconectar de su trabajo por un tiempo.

Asistentes de terapia ocupacional y psicólogos: Green In House podría ser utilizado como una herramienta en la terapia ocupacional o en el tratamiento de trastornos de atención, enseñando a los pacientes a concentrarse en tareas, a ser conscientes del momento presente y a cuidar de algo más.

Estos son solo algunos ejemplos de cómo Green In House podría ser útil para los adultos. Este producto tiene un amplio alcance debido a su combinación única de tecnología y educación práctica en el cuidado de las plantas.






Desglose detallado de cada objetivo de la introduccion

tambien objetivos personales de que quiero conseguir con el proyecto


Este apartado explica de forma precisa y concisa cuales son los objetivos que se persiguen con la realización del proyecto. Se puede distinguir entre los objetivos marcados por los requisitos del software a construir y los objetivos de carácter técnico que plantea a la hora de llevar a la práctica el proyecto.

\capitulo{3}{Conceptos teóricos}

Green In House es un proyecto que combina diversas tecnologías para facilitar el cuidado de las plantas, por lo que es necesario tener claro tanto los conceptos teóricos referentes a los factores que influyen en el crecimiento de las plantas, como los conceptos teóricos de las diferentes tecnologías involucradas en el sistema.

\section{Factores que influyen en el crecimiento de una planta}
Es esencial mantener un control adecuado de la temperatura, la humedad y la luminosidad para fomentar un desarrollo correcto de una planta.
\begin{itemize}
    \item \textbf{La temperatura} es un factor crítico en el crecimiento de las plantas. Las plantas, al igual que los humanos, tienen un rango de temperaturas máximas y mínimas en las que pueden sobrevivir y prosperar. Si la temperatura es demasiado baja, una planta puede entrar en un estado de latencia, en el cual el crecimiento se ralentiza o incluso se detiene, llegan a marchitar la planta. Por el contrario, si la temperatura es demasiado alta, puede dañar la estructura celular de la planta y causar estrés, lo cual puede resultar en un crecimiento pobre o en la marchitación de la planta. Por lo tanto, mantener la temperatura dentro del rango óptimo es crucial para garantizar un crecimiento saludable.
    \item \textbf{La humedad} juega un papel vital en la salud de las plantas. La humedad ambiental afecta a la tasa de transpiración. Este es el proceso mediante el cual el agua se evapora de las hojas de las plantas. Este proceso es esencial para el transporte de nutrientes a través de la planta, pero si la humedad es demasiado baja, la planta puede perder agua más rápidamente de lo que puede absorberla, llevándola a la deshidratación. Por otro lado, demasiada humedad puede promover el crecimiento de hongos y otras enfermedades. Por lo tanto, mantener el equilibrio correcto es crucial.  
    \item \textbf{La luminosidad} es fundamental para la fotosíntesis. Este es el proceso a través del cual las plantas convierten la luz solar en energía. Sin suficiente luz, una planta no puede producir la energía que necesita para crecer. Pero de nuevo, demasiada luz (especialmente la luz solar directa y fuerte), puede ser perjudicial y causar quemaduras en las hojas de la planta. Por eso es esencial garantizar que la planta reciba la cantidad adecuada de luz.
\end{itemize}
Además de estos tres factores en los que se centra Green In House, existen muchos más factores que influyen en el crecimiento de las plantas como por ejemplo la calidad del suelo, la cantidad de nutrientes que tiene la tierra, el abono que se utiliza, el espacio que dispone la planta para crecer, la necesidad de podar las ramas, etc. En esta memoria no se incide más sobre ellos ya que no serán controlados por Green In Houe, pero es relevante destacar la existencia de los mismos.

\section{Raspberry Pi: Microcomputador con entradas y salidas GPIO}
Raspberry Pi es una serie de computadoras de placa única, con  bajo coste y del tamaño de una tarjeta de crédito, la cual fue desarrollada en el Reino Unido por la Raspberry Pi Foundation, con el objetivo de promover la enseñanza de informática básica en las escuelas. Sin embargo, debido a su gran versatilidad y potencia de cómputo, se ha vuelto increíblemente popular en una variedad de aplicaciones, desde servidores web caseros hasta prototipos de productos electrónicos.

La Raspberry Pi ofrece varias ventajas frente a computadores y microcontroladores convencionales, al hacer una fusión de lo mejor de ambos mundos en una única placa. Algunas de las características que hacen destacar a esta placa son las siguientes:
\begin{itemize}
    \item \textbf{Interfaz de Entrada/Salida General (GPIO):} Uno de los aspectos más atractivos de la Raspberry Pi es su conjunto de pines de Entrada/Salida General (GPIO). Estos pines permiten la interacción con una amplia variedad de componentes electrónicos, como sensores y actuadores. Los pines GPIO de la Raspberry Pi son muy flexibles y pueden configurarse para leer o escribir datos digitales, utilizar buses de comunicación y en algunos modelos para leer y generar señales analógicas, lo que la convierte en una plataforma ideal para proyectos de automatización caseros.
    \item \textbf{Rendimiento y Almacenamiento:} A pesar de su pequeño tamaño, la Raspberry Pi es una computadora completamente funcional, la cual cuenta con procesadores de varios núcleos, permitiendo realizan tareas paralelas, y una tarjeta gráfica que permite desplegar un entorno gráfico y utilizarla como si de un ordenador convencional se tratase. Puede ejecutar varios sistemas operativos de carácter general como Windows y Linux, aunque dispone de su propia distribución Linux optimizada para su hardware, conocida como Raspbian. Esta es la distribución de sistema operativo que se ha elegido para Green In House. Además, dispone de un slot de tarjeta SD, el cual funciona como un disco duro, lo que permite tener un almacenamiento de datos local de gran tamaño, lo cual permite a Green In House no tener porblemas de espacio de almacenamiento de sus base de datos.
    \item \textbf{Red y Servidor:} La Raspberry Pi puede conectarse a las red a través de Ethernet o Wi-Fi (en los modelos que cuentan con esta característica). Esto permite a Green In House proporcionar y recoger datos a través de una API-REST. El modelo utilizado para el proyecto no dispone de conexión WiFi, por lo que ha sido necesario instalar un dongle WiFi y sus correspondientes controladores.
    \item \textbf{Costo y Comunidad:} La Raspberry Pi es una opción atractiva debido a su bajo costo (aunque actualmente debido a los problemas en la obtención de semiconductores y la inflacción, su precio se ha visto muy incrementado respecto a cuando salieron al mercado). A pesar de su potencia y flexibilidad, es significativamente más barata que la mayoría de las computadoras o microcontroladores especializados. Además, cuenta con una comunidad de usuarios activa y entusiasta que puede ser una excelente fuente de ayuda y recursos, al ser de caracter Open Source.
\end{itemize}

    \subsection{Raspbian: Sistema operativo de Raspberry Pi}
    Como se ha comentado anteriormente, existen diferentes versiones de sistema operativo soportados por la Raspberry Pi, o más bien dicho, modificados para poder funcionar en ella. 
    Para la realización de Green In House se ha optado por el uso de su distribución oficial Raspbian, en su versión 6.21, ya que está completamente diseñada para sacar el máximo partido a las características de la Raspberry Pi. Aunque existan otros SO optimizados que prometan dar mayor rendimiento, he decidido utilizar Raspbian en su versión estable, para evitar posibles problemas derivados de inconsistencias internas del sistema operartivo.

\section{Bash: Un Lenguaje de Programación de Shell para Linux}
Bash \citep{wiki:bash} es un intérprete de comandos de Unix y es el acrónimo de Bourne Again SHell. Se trata de una evolución del shell Bourne, que fue uno de los primeros intérpretes de comandos del sistema operativo Unix. Bash fue creado por Brian Fox y lanzado por primera vez en 1989. 

Bash es el shell por defecto en la mayoría de las distribuciones de Linux actuales y también se puede utilizar en otros sistemas operativos, como Windows y Mac, a través de software como CMD y Terminal. Bash se utiliza tanto para ejecutar comandos directamente en la terminal, como para crear scripts de shell, los cuales son programas que ejecutan una serie de comandos.

\begin{itemize}
    \item \textbf{Comandos:} Bash es un intérprete de comando, lo que significa que los usuarios pueden escribir comandos, los cuales Bash interpreta y ejecuta. Esto permite a los usuarios interactuar directamente con el sistema operativo y realizar tareas como gestionar archivos, iniciar y detener programas, configurar el entorno del sistema,...  
    \item \textbf{Scripts de Shell:} Bash también es un lenguaje de programación, lo que significa que los usuarios pueden escribir programas (denominado scripts) que ejecutan una serie de comandos. Estos scripts son especialmente útiles para automatizar tareas repetitivas. Un script de Bash puede incluir funciones, bucles, condicionales, y otras estructuras de control de flujo comunes de los lenguajes de programación.   
\end{itemize}
\subsection{Utilización de Bash en Green In House}
Green In House cuenta con varios scripts Bash desarrollados para facilitar las tareas de instalación, despliegue, ejecución y parada de la aplicación.

\section{Python: Lenguaje de Programación de Alto Nivel para desarrollo de aplicaciones de carácter general}
Python \citep{wiki:python} es un lenguaje de programación interpretado, de alto nivel y de propósito general. Fue creado por Guido van Rossum y lanzado por primera vez en 1991. Algunas de sus características principales son:
\begin{itemize}
    \item \textbf{Legibilidad y Sintaxis Clara:} Python fue diseñado con una gran énfasis en la legibilidad y la simplicidad, lo que facilita el aprendizaje del lenguaje. Las estructuras de control de Python, las funciones, los bucles y las sentencias condicionales, están diseñadas para ser fácilmente comprensibles.  
    \item \textbf{Interpretado:} Python es un lenguaje interpretado, lo que significa que el código Python se ejecuta línea por línea. Esto facilita la depuración de los programas y la experimentación interactiva en la terminal de Python o en entornos como Jupyter Notebook.   
    \item \textbf{Multiparadigma:} Python admite varios paradigmas de programación, incluyendo programación orientada a objetos, programación imperativa y programación funcional, lo que le da una gran flexibilidad para resolver problemas de diferentes maneras.  
    \item \textbf{Ecosistema de Bibliotecas:} Python tiene un ecosistema de bibliotecas de terceros extremadamente rico, que lo hacen adecuado para una amplia gama de tareas. Bibliotecas como NumPy y SciPy para la computación científica, Django y Flask para el desarrollo web, y TensorFlow y PyTorch para el aprendizaje profundo, SQLalchemy para la gestión de base de datos, Open-API para el desarrollos de servidores Api-Rest, TKinter para el desarrollo de interfaces gráficas, ADAFruit para el uso de sensores electrónicos,... hacen de Python una opción popular en muchos campos de la informática y la ciencia de datos.
\end{itemize}
Tal y como se describe en los puntos anteriores, Python es un lenguaje de programación dinámico y fuertemente tipado, que combina una sintaxis clara, una semántica simple y un ecosistema de bibliotecas de terceros amplio y robusto. Estas características han llevado a Python a ser uno de los lenguajes de programación más populares en el mundo. Además, debido a estas características, es uno de los lenguajes más utilizados para realizar prototipados y desarrollo ágil de proyectos, al permitir a los programadores realizar programas utilizando un número menor de líneas de código de las que se necesitarían escribir en otros lenguajes como C++ o Java. Debido a su extremada versatilidad, es muy apropiado para una amplia gama de tareas, desde el desarrollo de progragarmas simples, pasando por el desarrollo web, hasta el análisis de datos y la inteligencia artificial.
Todo el código funcional de Green In House está desarrollado en Python, apoyándose en el uso de librerías de terceros para la gestión de los diferentes servicios y tecnologías que utiliza e implementa.

\section{Flutter: Lenguaje de Programación de Alto Nivel para desarrollo de aplicaciones multiplataforma}
Flutter es un framework de desarrollo de software multiplataforma de código abierto creado por Google. Alagunas de sus carcterística más importantes son:
\begin{itemize}
    \item \textbf{Hot Reload:} Uno de los aspectos más útiles de Flutter es su capacidad de recarga en caliente. Esta característica permite a los desarrolladores experimentar, construir interfaces de usuario, agregar características y corregir errores más rápido. Los cambios en el código se reflejan en tiempo real sin necesidad de reiniciar la aplicación. 
    \item \textbf{Diseño Material y Cupertino:} Flutter está estrechamente integrado con el diseño Material de Google y el diseño Cupertino de Apple, permitiendo a los desarrolladores crear interfaces de usuario que siguen las directrices de diseño de cada plataforma.   
    \item \textbf{Rendimiento:} Flutter utiliza el lenguaje de programación Dart, que se compila en código nativo, lo que significa que las aplicaciones Flutter tienen un rendimiento cercano al nativo en todas las plataformas.    
    \item \textbf{Widgets:} Flutter utiliza un sistema de widgets para crear la interfaz de usuario. Los widgets son bloques de construcción modulares que incluyen elementos como botones, textos, barras de desplazamiento, conmutadores, etc.  
    \item \textbf{Soporte multiplataforma:} Flutter está diseñado para permitir a los desarrolladores crear aplicaciones de alto rendimiento y alta fidelidad para múltiples sistemas como iOS, Android y web, todo desde un mismo código
    \item \textbf{Integración de código nativo:} Flutter permite la integración de código nativo de los distintos sistemas que soporta, en caso de que se necesitase realizar alguna funcionalidad específica o crítica que requiera de desarrollo nativo. 
\end{itemize}
    \subsection{Dart: El lenguaje de programación de Flutter}
    El lenguaje de programación utilizado en Flutter es Dart, el cual también ha sido desarrollado por Google. Dart es un lenguaje orientado a objetos que se compila "Just in Time" (JIT) para el desarrollo y "Ahead of Time" (AOT) para la producción. Esto significa que durante el desarrollo, el código se puede compilar y cambiar sobre la marcha (JIT), lo que permite la recarga en caliente. Para la producción, el código se compila antes de tiempo (AOT) en código nativo, lo que resulta en un rendimiento muy eficiente.
    Dart es un lenguaje fuertemente tipapado, lo que significa que una vez que se ha creado una variable de un tipo, no es posible cambiarlo. Esto tiene sus desventajas, pero a la vez ayuda a evitar errores comunes en tiempo de ejecución, los cuales podrían hacer caer la aplicación. A pesar de su seguridad de tipos, Dart sigue siendo un lenguaje fácil de aprender y usar, especialmente para aquellos que ya están familiarizados con otros lenguajes de programación orientados a objetos, como Java, C++ o Python.
    \subsection{Utilización de Flutter en Green In House}
    Flutter ha sido utilizado para desarrollar una aplicón mobil capaz de utilizar los endpoints de la Api Rest integrada en Green In House para recoger los datos de la base de datos alojada en la Raspberry y poder trabajar con ellos y graficarles, para mostrar al usuario información de manera intuitiva.

\section{Git: Sistema de control de versiones}
Git es un sistema de control de versiones distribuido de código abierto que permite gestionar y ver los cambios que se han ido realizando un proyecto a lo largo del tiempo. Está diseñado para manejar todo tipo de proyectos (ya sean pequeños o muy grandes), con una gran velocidad y eficiencia.Algunas de las características más destacadas de Git son:
\begin{itemize}
    \item \textbf{Sistema distribuido:} La característica más importante de Git es que es un sistema de control de versiones distribuido. Esto significa que cada desarrollador tiene una copia completa del historial de cambios del proyecto en su máquina local. Esta característica permite trabajar de manera descentralizada y autónoma, sin tener que estar constantemente sincronizado con un servidor central.
    \item \textbf{Integridad de los datos:} Git está diseñado con un fuerte énfasis en la integridad de los datos. Cada cambio o "commit" en Git tiene una suma de comprobación asociada que se utiliza para verificar la integridad de los datos. Esto permite a los desarrolladores asegurarse de que, una vez que se ha registrado un cambio en Git, los archivos no se alterarán sin su conocimiento.
    \item \textbf{Ramas:} Git ofrece un soporte robusto y flexible para la ramificación y fusión de código. Esta característica permite a los desarrolladores trabajar en características y experimentos de manera aislada, sin afectar el código principal del proyecto. Una vez que una característica está lista, Git proporciona herramientas para fusionar ese trabajo con la rama principal.
    \item \textbf{Desempeño:} Git ha sido diseñado con un fuerte enfoque en el rendimiento. A pesar de que maneja historiales de cambios complejos y grandes cantidades de archivos y cambios, las operaciones en Git son generalmente muy rápidas. Esta es una característica de gran valor ya que permite realizar un desarrollo ágil y eficiente.
    \item \textbf{Plataforma de colaboración:} Combinado con plataformas de alojamiento de código como GitHub, Git se convierte en una poderosa plataforma de colaboración, permitiendo a los equipos de desarrolladores revisar y discutir cambios, gestionar tareas e integrar de manera continua el código. Este sistema de colaboración es esencial para mantener la calidad del código y coordinar el trabajo en equipo.
\end{itemize}
    \subsection{Utilización de Git en Green In House}
    Git ha sido utilizado en el desarrollo de Green In House junto con la plataforma GitHub para realizar el seguimiento de cambios, garantizar la integridad del código, detección de problemas tras la realización de cambios y poder mantener en la nube una copia de seguridad del proyecto actualizada continuamente en caso de que se estropease el entorno de desarrollo, evitando la pérdida del código del proyecto.




\section{Seccion}

En aquellos proyectos que necesiten para su comprensión y desarrollo de unos conceptos teóricos de una determinada materia o de un determinado dominio de conocimiento, debe existir un apartado que sintetice dichos conceptos.

Algunos conceptos teóricos de \LaTeX \footnote{Créditos a los proyectos de Álvaro López Cantero: Configurador de Presupuestos y Roberto Izquierdo Amo: PLQuiz}.

\section{Secciones}

Las secciones se incluyen con el comando section.

\subsection{Subsecciones}

Además de secciones tenemos subsecciones.

\subsubsection{Subsubsecciones}

Y subsecciones. 


\section{Referencias}

Las referencias se incluyen en el texto usando cite \cite{wiki:latex}. Para citar webs, artículos o libros \cite{koza92}.


\section{Imágenes}

Se pueden incluir imágenes con los comandos standard de \LaTeX, pero esta plantilla dispone de comandos propios como por ejemplo el siguiente:

\imagen{escudoInfor}{Autómata para una expresión vacía}{.5}



\section{Listas de items}

Existen tres posibilidades:

\begin{itemize}
	\item primer item.
	\item segundo item.
\end{itemize}

\begin{enumerate}
	\item primer item.
	\item segundo item.
\end{enumerate}

\begin{description}
	\item[Primer item] más información sobre el primer item.
	\item[Segundo item] más información sobre el segundo item.
\end{description}
	
\begin{itemize}
\item 
\end{itemize}

\section{Tablas}

Igualmente se pueden usar los comandos específicos de \LaTeX o bien usar alguno de los comandos de la plantilla.

\tablaSmall{Herramientas y tecnologías utilizadas en cada parte del proyecto}{l c c c c}{herramientasportipodeuso}
{ \multicolumn{1}{l}{Herramientas} & App AngularJS & API REST & BD & Memoria \\}{ 
HTML5 & X & & &\\
CSS3 & X & & &\\
BOOTSTRAP & X & & &\\
JavaScript & X & & &\\
AngularJS & X & & &\\
Bower & X & & &\\
PHP & & X & &\\
Karma + Jasmine & X & & &\\
Slim framework & & X & &\\
Idiorm & & X & &\\
Composer & & X & &\\
JSON & X & X & &\\
PhpStorm & X & X & &\\
MySQL & & & X &\\
PhpMyAdmin & & & X &\\
Git + BitBucket & X & X & X & X\\
Mik\TeX{} & & & & X\\
\TeX{}Maker & & & & X\\
Astah & & & & X\\
Balsamiq Mockups & X & & &\\
VersionOne & X & X & X & X\\
} 

\capitulo{4}{Técnicas y herramientas}
A continuación se describen 

\section{Técnicas y metodologías}
Para poder realizar de manera óptima y eficiente el diseño de Green In House, con un carácter iterativo y con posibilidades de ampliar sus funcionalidades continuamente en cada iteración, se ha recurrido a la utilización de patrones de diseño y a una metodología ágil. 

    \subsection{Patrones de Diseño}
    Los patrones de diseño \cite{wiki:patrones_de_diseño} son soluciones probadas a problemas comunes que nos encontramos en el diseño de software. Son como plantillas que se pueden aplicar en diferentes situaciones, proporcionando un esquema que puede ayudar a resolver un problema en particular. Básicamente son una forma de reutilizar conocimientos y experiencias previas para hacer el código de una aplicación más eficiente, mantenible y comprensible.
    En Green In House se ha hecho uso de los siguientes patrones de diseño, los cuales están más detallados en el apartado de anexos:
    \begin{itemize}                
        \item Patrón de diseño \textbf{MVC} (Modelo-Vista-Controlador)
        \item Patrón de diseño \textbf{Adaptador}
        \item Patrón de diseño \textbf{Factoría}
        \item Patrón de diseño \textbf{Estrategia}
        \item Patrón de diseño \textbf{Plantilla}
        \item Patrón de diseño \textbf{Fachada}
        \item Patrón de diseño \textbf{Singleton}
    \end{itemize} 

    \subsection{Scrum: Metodología de gestión ágil}
    Scrum \cite{wiki:scrum} es un marco de trabajo de metodología ágil ampliamente utilizado en el desarrollo de software. Se basa en principios colaborativos, flexibilidad y entrega incremental. El objetivo principal de Scrum es permitir a los equipos de trabajo desarrollar productos de alta calidad de manera eficiente. Además permite realizar una entrega temprana y continua de valor al cliente, a la vez que se fomenta la mejora continua del proceso de desarrollo.
    Scrum se basa en un enfoque iterativo e incremental, en el que el trabajo se divide en ciclos continuos de desarrollo llamados \textit{sprints}. Cada sprint tiene una duración fija, generalmente de 1 a 4 semanas, y al final de cada sprint se entrega un incremento funcional del producto. Esta metodología permite ajustar y cambiar los requisitos a medida que se realizan los incrementos y se obtiene retroalimentación del cliente.
    Durante el \textit{sprint}, se llevan a cabo reuniones diarias llamadas \textit{Daily Scrum}, en las cuales el equipo comparte el progreso realizado, identifica posibles problemas y planifica las tareas para el día. También se realizan reuniones de planificación de \textit{sprint}, revisión de \textit{sprint} y retrospectivas, que ayudan al equipo a refinar y mejorar continuamente su trabajo.
    El equipo de desarrollo de Scrum se organiza de manera autónoma y autoorganizada. No hay jerarquías rígidas, sino que se fomenta la colaboración y la toma de decisiones conjunta. El equipo se compone de roles claramente definidos:
    \begin{itemize}
      \item \textbf{Product Owner:} es el responsable de representar los intereses del cliente y definir los requisitos del producto. También se encarga de mantener el \textit{product backlog}, priorizar las tareas y asegurar que se cumplan las expectativas del cliente.
      \item \textbf{Scrum Master:} es el facilitador del equipo Scrum. Su rol es asegurar que se sigan los principios y prácticas de Scrum, eliminar problemas y garantizar un ambiente de trabajo colaborativo y productivo.
      \item \textbf{Equipo de Desarrollo:} está formado por profesionales que llevan a cabo el trabajo de desarrollo del producto. Son multidisciplinares y autorganizados, asumiendo la responsabilidad de entregar un incremento de software funcional al final de cada \textit{sprint}.
    \end{itemize}
    \imagen{scrum}{Resumen de funcionamiento de metodología Scrum}{.7}

\section{Herramientas}
Green In House no es una simple maceta, si no que es un proyecto que aglutina una gran cantidad de funcionalidades y tecnologías diferentes, las cuales han sido desarrolladas utilizando las siguientes herramientas.

    \subsection{Venv: Entornos Virtuales en Python}
    Uno de los aspectos más críticos al desarrollar proyectos en Python es la gestión de dependencias. Con frecuencia, diferentes proyectos requieren diferentes versiones de las mismas librerías. La instalación global de dichas librerías se hace inviable, ya que puede generar conflictos, debido a que determinadas versiones pueden integrar métodos que otras no, o requerir otra sintaxis, etc. Debido a esto, los entornos virtuales de Python, o \textit{venv} desempeñan un papel crucial. \cite{wiki:venv}
    El módulo \textit{venv} de Python proporciona un soporte para crear entornos virtuales ligeros con sus propios directorios de sitio, aislados del intérprete Python global del sistema. Cada entorno virtual tiene su propio intérprete de Python, lo que permite instalar y gestionar paquetes de Python independientes para cada proyecto.

    \subsection{SQLAlchemy: librería de base de datos en Python}
    SQLAlchemy \cite{wiki:sqlalchemy} es una librería de Python que facilita la comunicación entre programas desarrollados en Python y las bases de datos SQL. Proporciona un conjunto completo de operaciones de persistencia de datos de alto nivel, así como un marco de abstracción de base de datos de SQL completo para aportar una mayor flexibilidad. 
    \begin{itemize}
        \item \textbf{ORM (Object Relational Mapper):} una de sus características más poderosas es su ORM (Object-Relational Mapping). Un ORM es una técnica de programación que facilita la conversión de datos entre sistemas incompatibles (en este caso, entre objetos Python y tablas SQL). En otras palabras, un ORM nos permite trabajar con bases de datos SQL en un lenguaje orientado a objetos, como Python. Esto significa que con el ORM de SQLAlchemy, se pueden crear clases en Python que se correspondan directamente con las tablas de la base de datos. Cada instancia de una clase representa una fila en la tabla correspondiente. Esta correspondencia entre objetos Python y tablas SQL nos permite interactuar con los datos de la base de datos de una manera más intuitiva y similar a como normalmente lo hace Python.
        \item \textbf{Abstracción de base de datos:} además del ORM, también proporciona un nivel de abstracción que nos permite cambiar entre diferentes sistemas de bases de datos SQL con pocos o ningún cambio en el código. Esto es útil si en un futuro se decidiese cambiar la base de datos subyacente de Green In Hosue o quisiese hacer el código más general para que funcione con diferentes sistemas de bases de datos.
    \end{itemize}

    \subsection{OpenAPI: librería de API REST en Python}
    OpenAPI \cite{wiki:openapi}, anteriormente conocido como Swagger, es una especificación para archivos de definición de API REST. Estos archivos de definición pueden estar escritos en YAML o JSON y describen los detalles de la API REST, entre los cuales por ejemplo se encuentran:
    \begin{itemize}
        \item Endpoints disponibles.
        \item Métodos HTTP soportados.
        \item Parámetros de las solicitudes.
        \item Respuestas esperadas.
        \item Modelos de los datos a comunicar.
    \end{itemize}
    La especificación OpenAPI se utiliza para documentar y describir APIs REST de forma estandarizada y fácil de entender, tanto para las personas, como para las máquinas. Esto facilita el desarrollo, la prueba, y la integración de servicios basados en APIs. Un beneficio adicional de la especificación OpenAPI es que puede generar automáticamente una interfaz de usuario interactiva, permitiendo a los usuarios o desarrolladores explorar y probar la API de manera intuitiva.
        \subsubsection{Microframework Connexion}
        Green In House utiliza el Microframework Connexion, el cual mapea el
        API especificado en el archivo \texttt{spec.yml} con OpenAPI en Python, facilitando la creación de los endpoints para que otras aplicaciones puedan conectar con el servicio, el cual por defecto es lanzado en el puerto 5000. Para poder acceder a estos endpoints hay que hacerlo utilizando de base la siguiente url: \texttt{http://192.168.1.240:5000/api/v1/}
        \subsubsection{JSON: Archivos de datos con formato específicos}
        JSON, que significa JavaScript Object Notation, es un formato ligero de intercambio de datos. JSON es un formato de texto que es completamente independiente del lenguaje pero utiliza convenciones que son familiares para los programadores de diversos lenguajes como C, C++, C\#, Java, JavaScript, Perl, Python, y muchos otros. Se utiliza como sistema de almacenaje y compartición de datos entre diferentes aplicaciones, al tener un formato independiente de la aplicación.
            \subsubsubsection{Formato de JSON}
            Un objeto JSON es un conjunto de pares clave-valor rodeado por llaves \{\}. Las claves son cadenas de caracteres y los valores pueden ser:
            \begin{itemize}
                \item Números.
                \item Cadenas de caracteres.
                \item Booleanos (true o false).
                \item Otros objetos JSON.
                \item Arrays de cualquiera de las opciones anteriormente. mencionadas. Un array JSON es una lista de valores rodeada por corchetes [].
            \end{itemize}

    \subsection{TKinter:  librería de interfaces gráficas en Python}
    TKinter \cite{wiki:tkinter} es una librería de Python utilizada para el desarrollo de interfaces gráficas de usuario (GUI). Es la interfaz estándar de Python para el kit de herramientas Tk, y es tanto simple como eficaz para la creación de una variedad de programas. TKinter proporciona un poderoso conjunto de widgets de interfaz de usuario:
    \begin{itemize}
        \item Botones.
        \item Cajetines de texto.
        \item Listas.
        \item Barras de desplazamiento.
        \item Etiquetas.
        \item Menús.
        \item Ventanas emergentes.
        \item Frames.
    \end{itemize}
    Todo esto permite crear interfaces de usuario altamente personalizables y funcionales.  
    \imagen{app_tkinter}{Captura de pantalla de la app gráfica de Green In House desarrollada en TKinter}{1}

    \subsection{Adafruit - CircuitPython: librería de sensores y actuadores en Python}
    Adafruit CircuitPython \cite{wiki:adafruit_circuit_python} es una implementación de Python diseñada para simplificar la experimentación con hardware de bajo coste. CircuitPython es una implementación de Python 3 diseñada para microcontroladores. Esta librería está desarrollada por Adafruit Industries, una empresa reconocida en el mundo de la electrónica y hardware de código abierto.    
    Una de las características clave de CircuitPython es su simplicidad. Proporciona una excelente plataforma para aquellos que están empezando a programar y trabajar con hardware. CircuitPython es capaz de interactuar directamente con el hardware. Esto incluye leer entradas y salidas digitales, utilizar la comunicación I2C y SPI, y controlar y leer sensores y actuadores. Esto permite a los desarrolladores y estudiantes explorar la interacción entre el software y el hardware de una manera intuitiva y fácil de entender.
    % \begin{itemize}
    %     \item \textbf{MCP3008:} conversor de 8 señales analógicas a digitalES (ADC).
    %     \item \textbf{DHT11:} sensor de humedad y temperatura ambiente de señal digital.
    %     \item \textbf{LM35:} sensor de temperatura ambiente de señal analógica.
    %     \item \textbf{FC28:} sensor de humedad de tierra de la maceta de señal analógica.
    %     \item \textbf{LDR:} sensor de luminosidad ambiente de señal analógica.
    %     \item \textbf{BH1750:} sensor de luminosidad ambiente de señal digital.  
    % \end{itemize}
    \begin{table}[ht]
    \resizebox{13cm}{!} {
    \centering
    \begin{tabular}{|l|c|c|c|c|c|}
    \hline \textbf{Sensor} & \textbf{Función} & \textbf{Tipo de medición} & \textbf{Zona} & \textbf{Tipo de Señal} \\
    \hline MCP3008 & Conversor & - & - & Analógico a Digital \\
    \hline DHT11 & Sensor & Humedad, Temperatura & Ambiente & Digital \\
    \hline LM35 & Sensor & Temperatura & Ambiente & Analógica \\
    \hline FC28 & Sensor & Humedad & Maceta & Analógica \\
    \hline LDR & Sensor & Luminosidad & Ambiente & Analógica \\
    \hline BH1750 & Sensor & Luminosidad & Ambiente & Digital \\
    \hline
    \end{tabular}
    }
    \caption{Modelos de sensores utilizados en Green In House}
    \end{table}

    \section{Git: sistema de control de versiones}
    Git \cite{wiki:git} es un sistema de control de versiones distribuido de código abierto que permite gestionar y ver los cambios que se han ido realizando un proyecto a lo largo del tiempo. Está diseñado para manejar todo tipo de proyectos (ya sean pequeños o muy grandes), con una gran velocidad y eficiencia.Algunas de las características más destacadas de Git son:
    \begin{itemize}
        \item \textbf{Sistema distribuido:} su característica más importante es que es un sistema de control de versiones distribuido. Esto significa que cada desarrollador tiene una copia completa del historial de cambios del proyecto en su máquina local. Esta característica permite trabajar de manera descentralizada y autónoma, sin tener que estar constantemente sincronizado con un servidor central.
        \item \textbf{Integridad de los datos:} está diseñado con un fuerte énfasis en la integridad de los datos. Cada cambio o \textit{commit} tiene una suma de comprobación asociada que se utiliza para verificar la integridad de los datos. Esto permite a los desarrolladores asegurarse de que, una vez que se ha registrado un cambio, los archivos no se alterarán sin su conocimiento.
        \item \textbf{Ramas:} ofrece un soporte robusto y flexible para la ramificación y fusión de código. Esta característica permite a los desarrolladores trabajar en características y experimentos de manera aislada, sin afectar el código principal del proyecto. Una vez que una característica está lista, proporciona herramientas para fusionar ese trabajo con la rama principal.
        \item \textbf{Desempeño:} ha sido diseñado con un fuerte enfoque en el rendimiento. A pesar de que maneja historiales de cambios complejos y grandes cantidades de archivos y cambios, las operaciones son generalmente muy rápidas. Esta es una característica de gran valor ya que permite realizar un desarrollo ágil y eficiente.
        \item \textbf{Plataforma de colaboración:} Combinado con plataformas de alojamiento de código como GitHub, se convierte en una poderosa plataforma de colaboración, permitiendo a los equipos de desarrolladores revisar y discutir cambios, gestionar tareas e integrar de manera continua el código. Este sistema de colaboración es esencial para mantener la calidad del código y coordinar el trabajo en equipo.
    \end{itemize}
        
    \subsection{GitHub: Plataforma \textit{open source} de repositorios}
    GitHub \cite{wiki:github} es una plataforma de desarrollo de software basada en la nube que utiliza el sistema de control de versiones Git. Con el paso del tiempo se ha convertido en una herramienta esencial para muchos desarrolladores y organizaciones debido a las capacidades que ofrece para colaborar y gestionar proyectos de software.
    GitHub permite a los desarrolladores almacenar sus proyectos en repositorios, que son esencialmente directorios de proyectos. Estos repositorios pueden ser públicos (cualquier persona puede ver y contribuir al código) o privados (restringe el acceso sólo a personas específicas).
    Una característica principal de GitHub es su enfoque en la colaboración. Los desarrolladores pueden hacer un \textit{fork} de un repositorio, lo que crea una copia del proyecto (de la versión seleccionada) en su propia cuenta de GitHub. Después de eso pueden hacer cambios en esa copia y, cuando estén listos, pueden enviar una \textit{pull request} al repositorio original. El propietario del repositorio original puede revisar estos cambios y optar por descartarlos o por fusionarlos con el proyecto principal.
    Github también integra un control de ramas dentro de un mismo repositorio. Una rama en Git es una referencia a uno de los \textit{commits}. Cuando creamos una rama, en realidad estamos creando una nueva línea de desarrollo en base al código existente en un determinado \textit{commit}. Esto significa que podemos realizar cambios en esa rama sin afectar a otras ramas. Esto permite a los desarrolladores trabajar en características específicas, pruebas o experimentos en aislamiento, lo que facilita la organización del código y reduce el riesgo de introducir errores en el código de producción.
    Las ramas son especialmente útiles en el desarrollo colaborativo, ya que cada colaborador puede trabajar en una rama separada sin interferir con el trabajo de los demás. Una vez que un colaborador ha terminado su trabajo en una rama, puede solicitar que se integre de nuevo en la rama principal (a menudo llamada \textit{master} o \textit{main}) a través de un \textit{pull request}.     
    Además de estas capacidades de colaboración y control de versiones, GitHub ofrece características como la gestión de incidencias para rastrear y resolver problemas, y acciones de GitHub para automatizar flujos de trabajo de desarrollo.
    \subsection{Utilización de GitHub en Green In House}
        Git ha sido utilizado en el desarrollo de Green In House junto con la plataforma GitHub para realizar el seguimiento de cambios, garantizar la integridad del código y detección de problemas tras la realización de cambios. También ha sido una herramienta útil para poder mantener en la nube una copia de seguridad del proyecto actualizada continuamente en caso de que se estropease el entorno de desarrollo, evitando la pérdida del código del proyecto.

    \subsection{VSCode: IDE de programación multilenguaje}
    VSCode \cite{wiki:vscode} es la abreviatura de Visual Studio Code, el cual es un editor de código fuente desarrollado por Microsoft. Es una herramienta extremadamente popular y versátil utilizada por desarrolladores de software en todo el mundo. Algunas de las características clave que tiene son:   
    \begin{itemize}
        \item \textbf{Editor de código:} proporciona un editor de código altamente personalizable con resaltado de sintaxis acorde al lenguaje utilizado, sugerencias inteligentes, sangría automática y otras características útiles para facilitar la escritura de código.        
        \item \textbf{Extensibilidad:} hay miles de extensiones disponibles que agregan funcionalidades adicionales, como soporte para diferentes lenguajes de programación, integración con sistemas de control de versiones, etc.        
        \item \textbf{Depuración:} incluye capacidades de depuración integradas que permiten ejecutar y depurar el código directamente desde el editor. Esto es especialmente útil para identificar y solucionar problemas de manera ágil en el código.        
        \item \textbf{Integración con Git:} tiene integración con el sistema de control de versiones de Git. Permite realizar de manera cómoda y ágil tareas comunes de Git, como confirmar cambios, crear ramas y fusionar código, etc.
        \item \textbf{Terminal integrada:} dispone de una terminal integrada que permite ejecutar comandos de línea de comandos directamente dentro del editor. Esto evita tener que alternar entre el editor y la terminal del sistema operativo, lo que mejora la eficiencia en el flujo de trabajo.
    \end{itemize}
    
    \subsection{Conexión SSH: Control remoto}
    SSH \cite{wiki:ssh} es la abreviatura de \textit{Secure Shell}. Es un protocolo de red que permite una comunicación segura entre dos sistemas a través de una red insegura. Proporciona un mecanismo de autenticación y cifrado robusto para proteger la información transmitida entre ambos sistemas. Algunos de los de usos más comunes de SSH son:
    \begin{itemize}
        \item \textbf{Acceso remoto}: permite acceder a un sistema remoto de forma segura. Esto es útil cuando necesitas administrar un servidor o realizar tareas en un sistema al que no tienes acceso físico.
        \item \textbf{Transferencia de archivos segura}: permite transferir archivos de manera segura entre sistemas utilizando herramientas como \textit{scp} o \textit{sftp}. Esto es útil cuando necesitas mover archivos de un sistema a otro de forma segura y confiable.
        \item \textbf{Túneles de red}: permite crear túneles de red seguros que redirigen el tráfico a través de una conexión cifrada. Esto se utiliza comúnmente para acceder de forma segura a servicios de red, como bases de datos o servidores web, a través de una red no confiable. 
        \item \textbf{Ejecución de comandos remotos}: permite ejecutar comandos en un sistema remoto sin tener que estar físicamente presente en él. Esto es especialmente útil cuando necesitas automatizar tareas o administrar múltiples sistemas de forma remota.
    \end{itemize}
    Para poder establecer una conexión SSH es necesario contar con los siguientes elementos:
    \begin{itemize}
        \item \textbf{Servidor SSH}: el sistema al que se desea acceder debe tener un servidor SSH instalado y configurado correctamente.
        \item \textbf{Cliente SSH}: es necesario tener instalado y configurado un cliente SSH en el sistema local desde el cual te conectarás al servidor SSH remoto. En sistemas basados en Unix, como Linux o macOS, el cliente SSH suele estar disponible de forma nativa. En sistemas Windows, se pueden utilizar programas como PuTTY, Git Bash o VSCode, que incluyen clientes SSH. 
        \item \textbf{Credenciales de acceso}: para poder establecer la comunicación entre cliente y servidor hay que contar con las credenciales de acceso adecuadas, como un nombre de usuario y una contraseña, o una clave SSH privada si se ha configurado.
    \end{itemize}
    
    \subsection{FlutterFlow: IDE de programación GUI para Flutter}
    Flutter Flow \cite{wiki:flutter_flow} es una plataforma online que permite diseñar y desarrollar aplicaciones móviles utilizando Flutter de manera visual y sin necesidad de escribir código. Proporciona una interfaz intuitiva y fácil de utilizar que te permite crear interfaces de usuario y lógica de aplicación de manera rápida y sencilla. Para los desarrolladores experimentados, es importante destacar que también permite el diseño de funcionalidades propias más complejas de las que inicialmente ofrece. Algunas de sus características principales son:
    \begin{itemize}
        \item \textbf{Diseño visual:} ofrece una forma sencilla de diseñar la interfaz de usuario de una aplicación de forma visual, arrastrando y soltando componentes en un lienzo. Esto te permite ver instantáneamente cómo se verá la aplicación mientras se va construyendo.
        \item \textbf{Sin necesidad de codificación:} una de las ventajas más destacadas de Flutter Flow de cara al público general es que no requiere conocimientos de programación para crear aplicaciones móviles. Permite construir la lógica de la aplicación utilizando una interfaz gráfica intuitiva en lugar de escribir código manualmente.
        \item \textbf{Generación automática de código Flutter}: aunque no es necesario escribir código, Flutter Flow genera automáticamente el código Flutter correspondiente a medida que se diseña la aplicación. Esto permite obtener un proyecto de Flutter completamente funcional, el cual puede ser personalizado y mejorado a mayores según las necesidades de la aplicación.    
        \item \textbf{Colaboración en tiempo real}: permite la colaboración en tiempo real, lo que significa que puedes trabajar en equipo en el diseño y desarrollo de una aplicación simultáneamente. Esto es especialmente útil cuando se trabaja junto con otros desarrolladores o diseñadores en un mismo proyecto.    
        \item \textbf{Integraciones y personalizaciones}: admite diversas integraciones y personalizaciones para ampliar las capacidades de la aplicación diseñada. Permite agregar servicios externos, como bases de datos o servicios en la nube, y personalizar la lógica de la aplicación para adaptarse a requisitos específicos.
    \end{itemize}
    En sus versiones de pago ofrece integración con el repositorio GitHUB, sistema de traducción automática de textos y despliegue en las tiendas de aplicaciones de Android e iOS. Para el desarrollo de Green In House he contado con una licencia especial de estudiante, la cual me permite utilizar las funciones anteriormente nombradas. Para solicitar dicha licencia tuve que solicitarla a Flutter Flow entregando los documentos necesarios que me acreditan como alumno universitario y en 48 horas me la concedieron.
\capitulo{5}{Aspectos relevantes del desarrollo del proyecto}

Este apartado pretende recoger los aspectos más interesantes del desarrollo del proyecto, comentados por los autores del mismo.
Debe incluir desde la exposición del ciclo de vida utilizado, hasta los detalles de mayor relevancia de las fases de análisis, diseño e implementación.
Se busca que no sea una mera operación de copiar y pegar diagramas y extractos del código fuente, sino que realmente se justifiquen los caminos de solución que se han tomado, especialmente aquellos que no sean triviales.
Puede ser el lugar más adecuado para documentar los aspectos más interesantes del diseño y de la implementación, con un mayor hincapié en aspectos tales como el tipo de arquitectura elegido, los índices de las tablas de la base de datos, normalización y desnormalización, distribución en ficheros3, reglas de negocio dentro de las bases de datos (EDVHV GH GDWRV DFWLYDV), aspectos de desarrollo relacionados con el WWW...
Este apartado, debe convertirse en el resumen de la experiencia práctica del proyecto, y por sí mismo justifica que la memoria se convierta en un documento útil, fuente de referencia para los autores, los tutores y futuros alumnos.

\capitulo{6}{Trabajos relacionados}

Actualmente en el mercado ya se ofrecen diversas alternativas a este producto, aunque todas tienen un enfoque más autónomo puramente dicho, lo cual en Green In House se ha eliminado esa parte de autosuficiencia, para involucrar lo máximo posible al usuario  en el desarrollo de su planta (en este caso el niño):

\textbf{Click and Grow}: es un jardín interior inteligente que permite cultivar plantas en casa. Viene con cápsulas de semillas que se insertan en el jardín inteligente para un crecimiento sin problemas. Los sensores y el sistema automatizado de riego aseguran que las plantas reciban la cantidad adecuada de agua, luz y nutrientes. Los precios varían desde alrededor de 100€ hasta 200€ dependiendo del tamaño del jardín inteligente.

\textbf{AeroGarden}: es otro sistema de jardín interior inteligente que permite cultivar plantas en interiores durante todo el año. También utiliza semillas en cápsulas y tiene un sistema de iluminación LED ajustable y un panel de control para recordar cuándo agregar agua y nutrientes. Los precios oscilan entre 100€ y 300€ dependiendo del modelo.

\textbf{Planty Square}: es un jardín modular inteligente que permite a los usuarios cultivar varias plantas a la vez. Cada módulo puede cultivar una planta y los módulos se pueden conectar entre sí para formar un jardín más grande. Cuenta con una aplicación que envía notificaciones para regar las plantas. El precio ronda los 100€ .

\textbf{SproutsIO}: otro sistema es un jardín interior de alta tecnología que permite cultivar plantas sin tierra. Utiliza la hidroponía y la aeroponía para cultivar plantas y tiene una aplicación que permite a los usuarios monitorear y controlar su jardín desde su teléfono. El precio es más elevado, alrededor de 800€.

Es importante aclarar que estos precios pueden variar por región debido a factores como los impuestos, los costos de envío y la tasa de cambio. Todos estos productos tienen en común la idea de utilizar la tecnología para facilitar el cultivo de plantas en interiores, pero cada uno tiene sus propias características únicas, al igual que Green In House. Para determinar un precio de venta para Green In House, habría que tener en cuenta el costo de producción, las características del producto, el precio de los competidores y la disposición a pagar del público, lo cual se detallará más en la parte de anexos.

\begin{table}[ht]
\resizebox{13cm}{!} {
\begin{tabular}{|l|l|l|l|l|l|}
\hline   & \textbf{Green In House} & \textbf{Click and Grow} & \textbf{AeroGarden} & \textbf{Planty Square} & \textbf{SproutsIO} \\ \hline
\textbf{Precio}   & Medio   & Medio/Alto   & Alto   & Medio/Alto   & Muy alto   \\ \hline
\textbf{Conectividad}   & Sí   & Sí   & Sí   & Sí   & Sí   \\ \hline
\textbf{Completamente automatizado}     & No   & Sí   & Sí   & Sí   & Sí   \\ \hline
\textbf{Educación Ambiental}   & Sí   & No   & No   & No   & No   \\ \hline
\textbf{Personalización}   & Sí   & No   & No   & Parcial   & No   \\ \hline
\textbf{Posibilidad de expansión}   & Sí    & No    & No    & No    & No    \\ \hline
\textbf{Open Source}   & Sí    & No    & No    & No    & No    \\ \hline
\end{tabular}
}
\caption{Comparativa de Green In House con la competencia}
\label{tab:comparativa competidores}
\end{table}


\capitulo{7}{Conclusiones y Líneas de trabajo futuras}

A continuación se exponen las conclusiones obtenidas del desarrollo de Green In House y las posibles lineas de trabajo futuras para continuar mejorando Green In House.

\section{Conclusiones}
Actualmente Green In House se encuentra en una fase de desarrollo muy avanzada y se puede considerar un producto mínimo viable. Los objetivos iniciales del proyecto han sido cubiertos completamente, desarrollando un sistema capaz de recoger datos de diversos sensores electrónicos reales, almacenarlos en una base de datos, procesarlos y enviarlos como consultas de una API REST y generar nuevas instancias de elementos del sistema en la base de datos desde peticiones de la API REST. Además se ha diseñado con una estructura modular fácil de ampliar, con capacidad de generar nuevos sensores y plantas dinámicamente a través de servicios de la API REST. También cuenta con una aplicación gráfica que permite al usuario introducir las credenciales su red WiFi doméstica para que Green In House pueda comunicarse a través de red con la aplicación móvil desarrollada para visualizar gráficos de los registros del sistema en diferente periodos de tiempo.

La realización de Green In House me ha permitido investigar diferentes tecnologías y metodologías actuales a un nivel de profundidad considerable, suficiente como para hacerme a la idea de las grandes posibilidades de diseño y desarrollo que otorgan. 
Aunque estoy especializado en la programación de sistemas de automatización industrial mediante PLCs y redes de comunicación industrial para envío de datos entre dispositivos, este ha sido mi primer proyecto de automatización realizado  en la plataforma Raspberry Pi. Esta primera aproximación la he realizado bastante en profundidad, teniendo que investigar cómo hacer uso de su interfaz de comunicación con sensores, sus lenguajes y protocolos soportados, cuales son las tensiones de trabajo que soporta y cómo realizar la conexión con los distintos dispositivos electrónicos que puede manejar.

Esta investigación me ha ayudado a descubrir que Raspberry Pi aporta un entorno de desarrollo ágil muy completo, al soportar Python para el desarrollo de sus aplicaciones. Además cuenta con numerosas librerías específicas de manejo de sensores aparte de las múltiples librerías soportadas por Python de carácter general para realizar tareas como despliegue de un gestor de base de datos ORM y servicios de comunicación HTML como API REST. Además, al ser un sistema con múltiples núcleos, permite la ejecución paralela de tareas mediante hilos, lo cuál es muy útil para poder realizar varias tareas simultáneamente. Raspberry Pi ha sido una gran elección para el desarrollo de Green In House, al tener todas las funcionalidades que necesitaba reunidas en una misma placa. Sin embargo, he podido descubrir, que no es el sistema más adecuado para realizar tareas críticas de lectura de sensores en tiempo real. Esto no ha afectado al proyecto ya que la lectura de los sensores de Green In House se realiza en un intervalo de minutos, pero si se realizase cada poco milisegundos, sería más óptimo utilizar Arduino para leer dichos sensores y después enviar esa información almacenada en paquetes a la Raspberry Pi.

La realización de este proyecto también me ha permitido investigar más en profundidad en el campo de las bases de datos, y los gestores ORM. Durante el grado había desarrollado alguna pequeña base de datos para algunos trabajos, pero siempre siguiendo algún esquema ya dado. Esta es la primera vez que he tenido que discernir por mi mismo cual es la mejor opción de diseño de datos para poder llevar a cabo el proyecto, cuál es la mejor forma de segmentar dichos datos en tablas individuales y cómo interrelacionarlas de manera óptima para tener todos los datos vinculados entre sí.

\section{Líneas de trabajo futuras}
La parte de \textit{backend} está perfectamente diseñada y es completamente funcional en el estado actual, tanto para generar nuevas plantas, tipos de plantas, sensores, consejos, asociar plantas y sensores, dar de baja cualquier elemento del sistema, modificar los datos de elementos del sistema, etc., pero se pueden realizar ciertas mejoras, sobre todo en la parte del \textit{frontend}:
\begin{itemize}
    \item Diseñar la parte de frontend del tal manera que pueda aprovechar todas las funciones que ofrece el \textit{backend}, de una manera intuitiva y llamativa para un niño. Por ello esta parte sería mejor desarrollarla en combinación con alguna persona que tenga un amplio conocimiento en el ámbito del trabajo con niños. 
    \item Mejorar la aplicación de la pantalla táctil incorporada en la maceta para permitir visualizar en tiempo real los datos de los sensores.
    \item Incluir comprobación de contraseña de red WiFi introducida por el usuario para verificar que cumple el estándar WPA, ya que si no lo cumple, almacenar esta credencial en el sistema puede generar un problema permanente en la conectividad de la red WiFi.
    \item En la parte de backend se podría desarrollar la comunicación con nuevos modelos de sensores para ampliar la cantidad de factores físicos que es capaz de medir.
    \item Incluir seguridad en el servidor API REST desarrollando un sistema de verificación de \textit{API KEY} o de \textit{token}.
    \item Generar un sistema de notificaciones que avisase al usuario automáticamente cuando los valores de los sensores se saliesen de los rangos establecidos en los consejos de la planta.
    \item Intentar resolver en Flutter la dirección IP en la que está alojado el servidor API REST a través del hostname GreenInHouse, para poder habilitar el direccionamiento por DHCP y evitar posibles conflictos de red.
    \item Al poder generarse nuevos consejos para las plantas y modificar los existentes, también sería interesante crear una comunidad de usuarios donde pudieran compartir sus experiencias y consejos, para intentar aprender de otras personas. Esto sería muy interesante, ya que va acorde con los objetivos sociales de Green In House y promueve una interacción activa con otras personas.
    \item Generar un plan de pruebas automatizado de pruebas unitarias y de pruebas de integración para poder validar más eficientemente que al añadir un nuevo módulo o funcionalidad, no se ve repercutido el correcto funcionamiento del código existente.
    \item Otra posible mejorar que se podría realizar sería dockerizar la aplicación de Green In House para hacer aún más sencillo su despliegue en nuevos entornos.
    \item Otro aspecto en el que se podría ampliar la funcionalidad de Green In House es generar otro modelo totalmente automatizado e incluir actuadores que permitan realizar el riego automáticos o dar luz a la planta. Esto desviaría el proyecto de su principal enfoque educativo y su intención de obligar a los niños a interactuar los máximo posible con la planta, para aprender de este proceso. Por ello, si se quiere desarrollar, lo mejor sería desarrollar otro modelo independiente enfocado al cultivo automatizado.
\end{itemize}



\bibliographystyle{plain}
\bibliography{bibliografia}

\end{document}
